\pdfbookmark[0]{English title page}{label:titlepage_en}
\aautitlepage{%
  \englishprojectinfo{
  3. %Semester nr.
  }{
  \reporttitle
  }{%
    \subject
  }{%
    Fall Semester \the\year %project period
  }{%
    9 % project group
  }{%
    \authors
  }{%
    %list of supervisors
    Sokol Kosta
  }{%
    \today % date of completion
  }%
}{%department and address
    \textbf{Aalborg University Copenhagen} \\
    A.C. Meyers Vænge 15\\
    2450 København SV \\
    \href{http://www.aau.dk}{http://www.aau.dk}\\
    \textbf{Semester Coordinator:}  Hossam Mohamed Farag\\
    \textbf{Secretary:}  Nele Pernille Staun Hvid
}{
    % the abstract
The project focuses on developing the Lapsus solution for elderly citizens, supporting the UN’s Sustainable Development Goal 3.8 by ensuring easy access to quality healthcare through fall detection. Lapsus consists of a fall-detecting wristband with an integrated emergency button and a app used by both the user and their relatives. The solution enables elderly users to live independently at home while ensuring help is available in case of falls or illness. An embedded system with a threshold-based, state-machine algorithm detects falls and notifies relatives, while a separate emergency button handles non-fall emergencies. The app also displays the user's real-time location and sends the phone’s location when an emergency is triggered. Furthermore, relatives always have access so view the user's real-time location in the app. Testing the algorithm proved challenging due to the difficulty of simulating realistic falls.
}
% Here are the typical kinds of information found in most abstracts:

% 1. the context or background information for your research; the general topic under study; the specific topic of your research
% 2. the central questions or statement of the problem your research addresses
% 3. what’s already known about this question, what previous research has done or shown
% 4. the main reason(s), the exigency, the rationale, the goals for your research—Why is it important to address these questions? Are you, for example, examining a new topic? Why is that topic worth examining? Are you filling a gap in previous research? Applying new methods to take a fresh look at existing ideas or data? Resolving a dispute within the literature in your field? . . .
% 5. your research and/or analytical methods
% 6. your main findings, results, or arguments
% 7. the significance or implications of your findings or arguments.